\section{Carbon Flow Analysis}
\label{sec:carbon-analysis}

\subsection{Problem Form and Runtime Scope}

This chapter documents the post-power-flow carbon tracing layer.  It does not
solve a new electrical operating point.  Instead, it maps a converged canonical
hybrid network solution into directed carrier flows, source intensities, and
bus-level carbon intensity metrics.

\subsection{Code Map}

The current implementation is concentrated in:
\begin{itemize}
  \item \texttt{src/analysis/carbon\_analysis.cpp} for graph construction, tracing, and matrix-based intensity calculations,
  \item canonical PF result structures from the power-flow layer for branch, converter, and storage transfers,
  \item shared component emission-factor fields from the information model.
\end{itemize}

\subsection{Pseudocode Map}

The code-aligned workflow reference is Algorithm A15 in
Section~\ref{sec:algorithms}.  This chapter then expands the algebra used in
that workflow.

\subsection{Current Status}

The carbon-analysis route is implemented as a reporting layer over the canonical
hybrid network.  The main remaining gaps are source normalization for some
higher-level hybrid overlays and richer explicit reporting for auxiliary
participation paths.

\subsection{Purpose and Required Inputs}

The carbon-analysis module is a post-power-flow reporting layer.  It does not
solve a new electrical operating point.  Instead, it consumes
\begin{itemize}
  \item a \texttt{HybridPowerSystem},
  \item a converged \texttt{PowerFlowResult},
  \item component emission-factor data already stored in the information model,
\end{itemize}
and produces source-to-load carbon tracing reports and matrix-based bus carbon
intensities on the canonical hybrid AC/DC network.

The public entry points are
\begin{itemize}
  \item \texttt{compute\_carbon\_analysis(sys, pf\_result, opt)},
  \item \texttt{compute\_carbon\_analysis(sys, pf\_opt, ca\_opt)},
\end{itemize}
where the second overload simply runs the main PF solver first and then invokes
the carbon-analysis routine.

\subsection{Implemented Scope and Current Boundary}

The implementation now operates on the same canonical projected network used by
the steady-state solver.  In particular, the carbon model includes
\begin{itemize}
  \item canonical AC buses and AC branches,
  \item canonical DC buses and DC branches,
  \item VSC AC/DC converters through \texttt{PowerFlowResult::vsc\_transfers},
  \item DCDC converters through \texttt{PowerFlowResult::dcdc\_transfers},
  \item energy routers through \texttt{PowerFlowResult::er\_port\_transfers},
  \item AC and DC storage as explicit discharge-side carbon sources and charge-side traced sinks.
\end{itemize}

The current boundary is still narrower than the full information-model catalog:
\begin{itemize}
  \item explicit carbon-source normalization now includes \texttt{Generator}, \texttt{StaticGenerator} on AC and DC, \texttt{RenewableGen}, \texttt{PVSystem}, and discharging \texttt{Storage},
  \item charging \texttt{Storage} is represented as an explicit traced sink, but \texttt{MobileStorage} still participates only through auxiliary slack-balance bookkeeping,
  \item VPP and microgrid exchange are still treated as auxiliary net-injection terms in slack reconstruction; they are not yet emitted as standalone carbon sources,
  \item if PF did not converge or both AC and DC bus sets are empty, the routine returns an empty result object.
\end{itemize}

\subsection{Canonical Hybrid Graph}

Let the projected AC and DC bus sets be
\begin{align}
\mathcal{N}^{ac},\qquad \mathcal{N}^{dc},\qquad
\mathcal{N} = \mathcal{N}^{ac} \cup \mathcal{N}^{dc}.
\end{align}
The carbon module builds a directed hybrid carrier graph on \(\mathcal{N}\).
Each directed carrier edge is represented as
\begin{align}
e = \left(u,v,F_e,R_e,L_e\right),
\end{align}
where
\begin{itemize}
  \item \(u\) is the sending node,
  \item \(v\) is the receiving node,
  \item \(F_e\) is the active power entering the carrier at the sending side,
  \item \(R_e\) is the active power received at the downstream side,
  \item \(L_e = F_e - R_e\) is the carrier loss.
\end{itemize}
The carrier set is
\begin{align}
\mathcal{E} = \mathcal{E}^{ac} \cup \mathcal{E}^{dc} \cup \mathcal{E}^{vsc} \cup \mathcal{E}^{dcdc} \cup \mathcal{E}^{er}.
\end{align}

\subsection{Source and Load Normalization}

\paragraph{Unified source model.}
Each explicit source is normalized to
\begin{align}
s = \left(\mathrm{id}_s,\ i_s,\ P_s,\ e_s\right),
\end{align}
where \(i_s \in \mathcal{N}\) is the attached hybrid node, \(P_s\) is the
actual active injection in MW, and \(e_s\) is the emission factor in
\(\mathrm{tCO_2/MWh}\).

The implemented source families are
\begin{align}
\mathcal{S} =
\mathcal{G}^{sync}
\cup \mathcal{G}^{static,ac}
\cup \mathcal{G}^{static,dc}
\cup \mathcal{G}^{ren}
\cup \mathcal{G}^{pv}
\cup \mathcal{G}^{stor,ac,dis}
\cup \mathcal{G}^{stor,dc,dis}.
\end{align}
The source emission factors are read from
\begin{align}
e_s \in
\left\{
\begin{aligned}
&\texttt{Generator::emission\_factor\_tco2\_mwh}, \\
&\texttt{StaticGenerator::co2\_emission\_rate}, \\
&0.0, \\
&\texttt{Storage::soc\_carbon\_intensity\_tco2\_mwh}
\end{aligned}
\right\},
\end{align}
so \texttt{RenewableGen} and \texttt{PVSystem} are currently treated as
zero-emission injections.

\paragraph{Storage as explicit source / sink.}
For storage unit \(s\), the PF sign convention is bus injection:
\begin{align}
p_s > 0 \Rightarrow \text{discharge / inject},\qquad
p_s < 0 \Rightarrow \text{charge / withdraw}.
\end{align}
The implemented carbon normalization is
\begin{align}
p_s > 0 &: \text{explicit source with } P_s = p_s,\quad e_s = w_s^{soc}, \\
p_s < 0 &: \text{explicit sink with } P_s^{ch} = -p_s,
\end{align}
where
\begin{align}
w_s^{soc} = \texttt{Storage::soc\_carbon\_intensity\_tco2\_mwh}.
\end{align}
Therefore discharge-side carbon leaves the storage at the current stored-energy
carbon intensity, while charging-side carbon is traced from the upstream source
mix and exported in \texttt{storage\_carbon}.

\paragraph{SOC carbon-intensity propagation rule.}
For a multi-period storage interpretation with interval \(\Delta t\), stored
energy \(E_{s,t}\), stored carbon inventory \(C_{s,t}\), self-discharge
\(\sigma_s\), charge efficiency \(\eta_s^{ch}\), discharge power
\(P_{s,t}^{dis}\), and charging power \(P_{s,t}^{ch}\), the recommended update
rule is
\begin{align}
E_{s,t+1} &= (1-\sigma_s) E_{s,t}
          + \eta_s^{ch} P_{s,t}^{ch}\Delta t
          - P_{s,t}^{dis}\Delta t, \\
C_{s,t+1} &= (1-\sigma_s) C_{s,t}
          + w_{s,t}^{in}\eta_s^{ch} P_{s,t}^{ch}\Delta t
          - w_{s,t}^{soc} P_{s,t}^{dis}\Delta t, \\
w_{s,t+1}^{soc} &= \frac{C_{s,t+1}}{E_{s,t+1}}.
\end{align}
The current C++ carbon-analysis module exports the instantaneous field
\texttt{soc\_carbon\_intensity\_tco2\_mwh} and uses it directly during discharge;
the dynamic update above is the governing rule for any future time-coupled
simulation or optimization layer.

\paragraph{Slack-source reconstruction.}
Non-slack synchronous generators keep their scheduled \texttt{pg\_mw}.  The
slack source is reconstructed from hybrid active-power balance:
\begin{align}
P_{slack} = \max\!\left(
\sum_{\ell \in \mathcal{L}^{ac,u} \cup \mathcal{L}^{dc,u} \cup \mathcal{L}^{stor,ch}} P_{\ell}
+ \sum_{e \in \mathcal{E}} L_e
+ \sum_{i \in \mathcal{N}^{ac}} G_i^{sh}
- \sum_{s \in \mathcal{S}_{ns}} P_s
- P^{aux}_{net},
0\right),
\end{align}
where \(G_i^{sh}\) is the active shunt demand on AC bus \(i\),
\(\mathcal{S}_{ns}\) is the explicit non-slack source set, and
\(P^{aux}_{net}\) is the residual net injection from auxiliary components not
yet promoted to explicit carbon sources or sinks, currently mobile storage, VPP,
and microgrid exchange terms.

\paragraph{Unified load model.}
The module constructs
\begin{align}
\mathcal{L}^{ac,u},\qquad \mathcal{L}^{dc,u},\qquad \mathcal{L}^{stor,ch},
\end{align}
by combining first-class load rows with synthetic bus-demand rows created from
\texttt{ACBus::pd\_mw} or \texttt{DCBus::pd\_mw} when no explicit component row
covers that bus, plus one charge-side sink row per in-service storage unit with
\(p_s<0\).  Each unified load is represented as
\begin{align}
\ell = \left(\mathrm{id}_{\ell},\ i_{\ell},\ P_{\ell},\ Q_{\ell}\right).
\end{align}

\subsection{Directed Carrier Construction}

\paragraph{AC branches.}
For AC branch \(k=(i,j)\), let the solved terminal injections be
\(P_{f,k}\) and \(P_{t,k}\).  The branch loss is
\begin{align}
L_k = P_{f,k} + P_{t,k}.
\end{align}
The directed-edge rule is
\begin{align}
P_{f,k} > 0 &:\ i \rightarrow j,\quad F_k = P_{f,k},\quad R_k = \max(-P_{t,k}, 0), \\
P_{t,k} > 0 &:\ j \rightarrow i,\quad F_k = P_{t,k},\quad R_k = \max(-P_{f,k}, 0).
\end{align}
Zero-flow AC branches are omitted from the directed carrier graph.

\paragraph{DC branches.}
For DC branch \(k=(i,j)\) with resistance \(r_k\), solved endpoint voltages
\(V_i\) and \(V_j\), and system base \(S_{base}\), the implementation
reconstructs the DC current and terminal powers as
\begin{align}
I_k &= \frac{V_i - V_j}{r_k}, \\
P_{i\rightarrow k} &= S_{base} V_i I_k, \\
P_{j\rightarrow k} &= - S_{base} V_j I_k.
\end{align}
If \(P_{i\rightarrow k} > 0\), then the edge direction is \(i \rightarrow j\) with
\begin{align}
F_k = P_{i\rightarrow k},\qquad
R_k = \max(-P_{j\rightarrow k}, 0),\qquad
L_k = P_{i\rightarrow k} + P_{j\rightarrow k}.
\end{align}
The opposite direction is used when \(P_{j\rightarrow k} > 0\).

\paragraph{VSC converters.}
The PF sign convention is bus injection on both sides:
\begin{align}
p^{ac}_c + p^{dc}_c + L_c = 0.
\end{align}
If the AC side withdraws power from the AC bus,
\begin{align}
p^{ac}_c < 0,\qquad p^{dc}_c > 0,
\end{align}
then the converter transports carbon from AC to DC with
\begin{align}
F_c = -p^{ac}_c,\qquad R_c = p^{dc}_c,\qquad L_c = -(p^{ac}_c + p^{dc}_c).
\end{align}
If instead \(p^{dc}_c < 0\) and \(p^{ac}_c > 0\), the same rule is applied in
DC-to-AC direction.

\paragraph{DCDC converters.}
For DCDC transfer \((p^{in}_d, p^{out}_d, L_d)\), the bus injections are
\begin{align}
P^{inj}_{bus\_in} = -p^{in}_d,\qquad
P^{inj}_{bus\_out} = p^{out}_d,\qquad
L_d = p^{in}_d - p^{out}_d.
\end{align}
When \(p^{in}_d > 0\) and \(p^{out}_d > 0\), the carbon carrier direction is
\(bus\_in \rightarrow bus\_out\) with
\begin{align}
F_d = p^{in}_d,\qquad R_d = p^{out}_d.
\end{align}
If the signs reverse, the directed carrier is flipped accordingly.

\paragraph{Energy routers.}
Each router port uses the same bus-injection convention:
\begin{align}
p_{r,p} < 0 \Rightarrow \text{bus} \rightarrow \text{router input},\qquad
p_{r,p} > 0 \Rightarrow \text{router} \rightarrow \text{bus output}.
\end{align}
For router \(r\), define the active input and output port sets
\begin{align}
\mathcal{P}^{in}_r = \{p \mid p_{r,p} < 0\},\qquad
\mathcal{P}^{out}_r = \{p \mid p_{r,p} > 0\},
\end{align}
and the aggregate powers
\begin{align}
P_r^{in} = \sum_{p \in \mathcal{P}^{in}_r} (-p_{r,p}),\qquad
P_r^{out} = \sum_{p \in \mathcal{P}^{out}_r} p_{r,p},\qquad
L_r = \max(P_r^{in} - P_r^{out}, 0).
\end{align}
The carbon module reconstructs a proportional internal carrier graph between
all input/output port pairs.  For input port \(i\) and output port \(o\),
\begin{align}
\beta_o &= \frac{p_{r,o}}{P_r^{out}}, \\
F_{io} &= (-p_{r,i}) \beta_o, \\
R_{io} &= F_{io} \cdot \min\!\left(\frac{P_r^{out}}{P_r^{in}}, 1\right), \\
L_{io} &= F_{io} - R_{io}.
\end{align}
This construction preserves the solved port-sign convention, allocates router
losses from actual port mismatch, and avoids inventing additional loss from the
metadata field \texttt{loss\_percent}.

\subsection{Proportional Tracing on the Hybrid Graph}

\paragraph{Node inflow composition.}
For each hybrid node \(n \in \mathcal{N}\), the total inflow used for tracing is
\begin{align}
I_n = \sum_{s \in \mathcal{S}_n} P_s + \sum_{e \in \delta^{-}(n)} R_e,
\end{align}
where \(\mathcal{S}_n\) is the set of explicit sources attached to node \(n\).

\paragraph{Load tracing.}
For each load \(\ell\) attached to node \(n\), the initial traced fraction is
\begin{align}
\phi_{\ell,n}^{(0)} = \frac{P_{\ell}}{I_n}.
\end{align}
During the upstream breadth-first traversal, every visited node \(m\) allocates
its local explicit sources by
\begin{align}
C_{\ell,s} \mathrel{+}= \phi_{\ell,m} P_s,
\qquad s \in \mathcal{S}_m,
\end{align}
and each incoming carrier edge \(e:u\rightarrow m\) propagates the traced
fraction in proportion to the received-flow share
\begin{align}
\phi_{\ell,u} \leftarrow \phi_{\ell,m} \frac{R_e}{I_m}.
\end{align}
The resulting source-attribution map \(C_{\ell,s}\) is exported as
\texttt{LoadCarbonResult::generator\_supply\_mw} for both AC and DC load
reports.

\paragraph{Carrier-loss tracing.}
For each lossy carrier edge \(e\), the implementation traces the source
composition at the sending node and allocates the loss proportionally:
\begin{align}
\Lambda_{e,s} \approx \gamma_{e,s} L_e,
\end{align}
where \(\gamma_{e,s}\) is the traced source share at the carrier input.  These
allocations populate
\begin{itemize}
  \item \texttt{branch\_carbon} for AC branches,
  \item \texttt{dc\_branch\_carbon} for DC branches,
  \item \texttt{vsc\_carbon} for VSC losses,
  \item \texttt{dcdc\_carbon} for DCDC losses,
  \item \texttt{energy\_router\_carbon} for router internal losses,
\end{itemize}
as well as the aggregated \texttt{generator\_loss\_allocation} map.

\paragraph{Tracing-derived carbon intensity.}
For any traced demand or lossy carrier item \(x\), the weighted-average carbon
intensity is
\begin{align}
w_x^{trace} = \frac{\sum_s e_s C_{x,s}}{\sum_s C_{x,s}},
\end{align}
and the reported emissions are
\begin{align}
E_x^{trace} = w_x^{trace} P_x,
\end{align}
with \(P_x\) equal to demand MW for loads or loss MW for carriers.

\subsection{Matrix-Based Hybrid Bus Carbon Intensity}

The second path solves one carbon-intensity value per hybrid node.  Define
\begin{align}
w = \begin{bmatrix} w^{ac} \\ w^{dc} \end{bmatrix},
\end{align}
and let \(\alpha \in [0,1]\) control how carrier losses are split between the
sending and receiving sides.

\paragraph{Node-level bookkeeping.}
For each directed carrier edge \(e:u\rightarrow v\), the implementation applies
\begin{align}
P_{in,v} &\mathrel{+}= F_e - \alpha L_e, \\
P_{out,u} &\mathrel{+}= F_e, \\
P_{out,v} &\mathrel{+}= (1-\alpha)L_e.
\end{align}
Explicit source injections are added to \(P_{in}\), while unified AC/DC loads
and AC shunt conductance demand are added to \(P_{out}\).

\paragraph{Linear system.}
The upstream transfer matrix is
\begin{align}
A_{in}(i,j) =
\begin{cases}
F_e - \alpha L_e, & \text{if there is a directed carrier } j \rightarrow i, \\
0, & \text{otherwise},
\end{cases}
\end{align}
and the hybrid system matrix is
\begin{align}
A = \mathrm{diag}(P_{out}) - A_{in},
\qquad
b_i = \sum_{s \in \mathcal{S}_i} e_s P_s.
\end{align}
The node carbon-intensity vector is obtained from
\begin{align}
A w = b.
\end{align}

\paragraph{End-node handling.}
When a node has neither local generation nor local load and its system-matrix
row is nearly singular, the implementation replaces that row with an averaging
constraint
\begin{align}
w_i = \frac{1}{|\mathcal{N}(i)|} \sum_{j \in \mathcal{N}(i)} w_j,
\end{align}
so electrically passive end-nodes inherit the intensity of adjacent hybrid
nodes.  The current linear solver then converts the repaired sparse system to a
dense matrix and solves it with column-pivoted Householder QR.  The reported
\texttt{matrix\_residual} is
\begin{align}
\|Aw-b\|_2.
\end{align}

\subsection{Emissions Summaries}

The tracing path reports a generation-vs-load-vs-loss balance:
\begin{align}
E_{gen} &= \sum_{s \in \mathcal{S}} e_s P_s, \\
E_{load}^{trace} &=
\sum_{\ell \in \mathcal{L}^{ac,u}} E_{\ell}^{trace}
+ \sum_{\ell \in \mathcal{L}^{dc,u}} E_{\ell}^{trace}
+ \sum_{\ell \in \mathcal{L}^{stor,ch}} E_{\ell}^{trace}, \\
E_{loss}^{trace} &= \sum_{e \in \mathcal{E}} E_e^{trace}.
\end{align}
The balance error is
\begin{align}
\Delta E^{trace} = E_{gen} - E_{load}^{trace} - E_{loss}^{trace},
\end{align}
and the verification flag \texttt{tracing\_verified} checks whether the relative
balance error is below \texttt{verification\_tol}.

The matrix path reports an analogous summary using solved hybrid node
intensities:
\begin{align}
E_{load}^{mat} &= \sum_{\ell \in \mathcal{L}^{ac,u} \cup \mathcal{L}^{dc,u} \cup \mathcal{L}^{stor,ch}}
                 w_{i(\ell)} P_{\ell}, \\
E_{loss}^{mat} &= \sum_{e=(u,v) \in \mathcal{E}} w_u L_e.
\end{align}
The matrix loss summary therefore uses the sending-node intensity of each
carrier rather than the average of the two terminal nodes.

\subsection{Result Artifacts}

The dedicated carbon result object now contains
\begin{itemize}
  \item \texttt{load\_carbon} and \texttt{dc\_load\_carbon}: per-load tracing allocations and carbon intensity,
  \item \texttt{branch\_carbon} and \texttt{dc\_branch\_carbon}: per-branch loss allocations and carbon intensity,
  \item \texttt{vsc\_carbon} and \texttt{dcdc\_carbon}: converter-loss allocations and hybrid transfer direction,
  \item \texttt{storage\_carbon}: per-storage charge/discharge carbon summary, including traced charging intensity and explicit discharge-side SOC intensity,
  \item \texttt{energy\_router\_carbon}: per-router internal-loss carbon allocation reconstructed from port transfers,
  \item \texttt{bus\_carbon} and \texttt{dc\_bus\_carbon}: matrix-based AC/DC node carbon intensities,
  \item \texttt{generator\_loss\_allocation}: aggregated traced loss contribution by source id,
  \item \texttt{tracing\_summary} and \texttt{matrix\_summary},
  \item \texttt{tracing\_verified}, \texttt{matrix\_solved}, and \texttt{matrix\_residual}.
\end{itemize}

The field-level JSON / workbook contracts for these artifacts are documented in
the I/O-contract and output-dictionary sections of this notebook; this chapter
focuses on the analytical model and the meaning of the reported quantities.

\subsection{Sankey Visualization of Carbon/SO\texorpdfstring{$_2$}{2} Source-Sink Relations}

For engineering review, the repository includes a lightweight post-processing
workflow to render source-to-sink Sankey charts:
\begin{enumerate}
  \item run PF + carbon analysis and export a carbon-result JSON using \texttt{run\_carbon\_analysis\_cli},
  \item render Sankey HTML using \texttt{tools/plot\_emission\_sankey.py}.
\end{enumerate}

The standard commands are:
\begin{verbatim}
cmake --build build --target run_carbon_analysis_cli -j4
./build/run_carbon_analysis_cli \
  --input-system-json <system.json> \
  --output-carbon-json <carbon_result.json>

python3 tools/plot_emission_sankey.py \
  --system-json <system.json> \
  --carbon-json <carbon_result.json> \
  --output-html <emission_sankey.html> \
  --aggregate-view bus \
  --top-n-source-nodes 20 \
  --top-n-sink-nodes 30 \
  --top-n-links 200 \
  --load-granularity bus \
  --loss-granularity type
\end{verbatim}

By default, the generated HTML contains two Sankey graphs:
\begin{itemize}
  \item carbon flow (\(\mathrm{tCO_2/h}\)),
  \item sulfur-dioxide flow (\(\mathrm{kg/h}\)).
\end{itemize}
The SO\(_2\) graph currently reads source factors from
\texttt{Generator::so2\_factor\_kg\_mwh}; non-generator source families are
treated as zero-SO\(_2\) unless optional extension keys are provided in the
input JSON.
